\documentclass[11pt, a4paper]{article}

% --- Packages ---
\usepackage{amsmath, amssymb, amsthm}
\usepackage{mathtools}
\usepackage{hyperref}
\usepackage{geometry}
\usepackage{graphicx}
\usepackage{booktabs}
\usepackage{cleveref}
\usepackage[style=authoryear, backend=biber]{biblatex}
\addbibresource{refs.bib}

\geometry{margin=2.5cm}

% --- Theorem environments ---
\newtheorem{theorem}{Theorem}
\newtheorem{proposition}[theorem]{Proposition}
\newtheorem{lemma}[theorem]{Lemma}
\newtheorem{corollary}[theorem]{Corollary}
\newtheorem{definition}{Definition}
\newtheorem{remark}{Remark}

% --- Notation ---
\newcommand{\E}{\mathbb{E}}
\newcommand{\Q}{\mathbb{Q}}
\newcommand{\R}{\mathbb{R}}
\newcommand{\dd}{\,\mathrm{d}}
\newcommand{\ii}{\mathrm{i}}   % imaginary unit

\title{\textbf{The Heston Stochastic Volatility Model}\\
       \large Semi-Analytical Pricing, Monte Carlo Simulation, and Market Calibration}
\author{Lucas Hutton}
\date{\today}

\begin{document}
\maketitle
\tableofcontents
\newpage

% ===========================================================================
\section{Introduction}
% ===========================================================================
% TODO: motivate stochastic volatility — the vol smile, limitations of BS.
% Introduce the Heston model and what this write-up covers.


% ===========================================================================
\section{Mathematical Preliminaries}
% ===========================================================================

\subsection{Brownian Motion and It\^{o} Calculus}
% TODO: define Brownian motion, state Itô's lemma (one and multi-dimensional).

\subsection{Risk-Neutral Pricing and Girsanov's Theorem}
% TODO: change of measure, fundamental theorem of asset pricing, the
% risk-neutral expectation formula  C = e^{-rT} E^Q[payoff].


% ===========================================================================
\section{The Black--Scholes Model}
% ===========================================================================
% TODO: GBM, BS PDE, closed-form call formula, implied volatility, vol smile.
% This section exists to motivate WHY we need stochastic vol.



\end{document}